\section{Results}
\label{sec:results}

\subsection{Comparison with Existing Systems}
\label{sec:results-main}
% ============================================================

\paragraph{Retrieval accuracy comparison with long-term memory systems
(Table~\ref{tab:memory-results}).}
Agents do not natively retain information across sessions, so answering questions over prior interactions requires an external memory system. 
We compare Scroll against dedicated long-term memory systems. For each system, we report the best publicly available result under its own preferred configuration (backbone model, retrieval budget, and judge) as of August~15, 2026.\footnote{We do not reproduce the baselines ourselves, as independent reproductions in this area have repeatedly led to disagreement over evaluation setup~\citep{zep-locomo-blog, mem0-zep-issue}.} 
% This convention favors the baselines, as each reported number constitutes an upper bound on that system's performance.
Appendix~\ref{app:breakdown} provides per-category breakdowns of Scroll on the \texttt{S} and \texttt{M} splits of LongMemEval and on BEAM$_{10M}$.

\begin{table}[t]
\centering
\caption{Comparison with existing long-term memory systems. For each baseline, we report the best publicly available result under that system's own setup as of August~15, 2026; ``–'' denotes no publicly reported result. These are reference points from the literature rather than a controlled comparison: reader models differ across rows and can substantially affect scores---EmergenceMem (GPT-4o), Zep (GPT-5.4), Mastra OM (GPT-5 mini), Mem0 (GPT-5), Hindsight (Gemini~3~Pro), Exabase M-1 (Gemini~3~Flash), RAG and LIGHT (Llama-4-Maverick); Honcho uses a multi-model pipeline and Cognee does not
report its reader.}
\label{tab:memory-results}
\small
\begin{tabular}{lcc}
\toprule
Method & LongMemEval$_S$ & BEAM$_{10M}$ \\
\midrule
RAG~\citep{beam}                     & --   & 24.9 \\
LIGHT~\citep{beam}                   & --   & 26.6 \\
\midrule
Zep~\citep{zep-research}             & 90.2 & --   \\
Mem0~\citep{mem0,mem0-eval}          & 94.4 & 48.6 \\
Hindsight~\citep{hindsight}          & 94.6 & 64.1 \\
EmergenceMem~\citep{emergencemem}    & 86.0 & --   \\
Honcho~\citep{honcho}                & 90.4 & 40.6 \\
Mastra OM~\citep{mastra}             & 94.9 & --   \\
Cognee~\citep{cognee-beam}           & --   & 67.0 \\
Exabase M-1~\citep{exabase}          & \textbf{96.4} & 68.0 \\
\midrule
Scroll (ours)                        & 94.8 & \textbf{73.1} \\
\bottomrule
\end{tabular}
\end{table}

As shown in Table \ref{tab:memory-results}, Scroll is competitive with the strongest reported systems on LongMemEval$_S$ and beats the best-performing system by 5.1 points on BEAM$_{10M}$. Existing memory systems follow a three-stage paradigm: at ingestion, an LLM processes the history into a derived store through fact extraction, summarization, or knowledge-graph construction; at query time, a retrieval pipeline selects candidate memories from that store; and a reader model then reasons over the returned snippets to produce the answer. Scroll instead ingests the raw history as-is, composes retrieval code per question, and needs no separate reader: the agent that writes and executes the queries also produces the final answer directly. 

\begin{table}[t]
\centering
\caption{LOCA accuracy (\%) of different context-management strategies at
the two largest environment description lengths. All agent loops use
Qwen3.8-Max as the backbone;
$\Delta$ denotes the absolute drop from 128K to 256K.}
\label{tab:loca-results}
\small
\begin{tabular}{lccc}
\toprule
Agent loop & 128K & 256K & $\Delta$ \\
\midrule
Summarization Agent            & 86.7 & 65.3 & -21.4 \\
Retrieval Agent           & 88.0 & 66.7 & -21.3 \\
CodeAct Agent                  & \textbf{89.3} & 85.3 & -4.0 \\
\midrule
Scroll (ours)                  & \textbf{89.3} & \textbf{86.7} & \textbf{-2.6} \\
\bottomrule
\end{tabular}
\end{table}

\paragraph{Comparison of context-management strategies on LOCA (Table~\ref{tab:loca-results}).}
On LOCA, we compare four agents that share the same backbone (Qwen3.8-Max) and toolset, and differ only in how they manage a growing context: (i) a \emph{summarization agent}, a ReAct agent~\citep{react} that periodically compacts its interaction history into a summary; (ii) a \emph{retrieval agent}, a ReAct agent whose overflowing history is evicted and made accessible through a recall tool; (iii) a \emph{CodeAct agent}~\citep{codeact} that interacts with the environment through programmatic tool calling; and (iv) Scroll. A comparison against the best published numbers from the LOCA paper~\citep{zeng2026loca} and its leaderboard is in Appendix~\ref{app:loca-full}.

Table~\ref{tab:loca-results} reports accuracy at the two largest environment description lengths. The CodeAct agent and the agent with Scroll, which bind intermediate results to environment objects instead of carrying raw text in context, achieve the best performance and the smallest decrease as context grows.


% ============================================================
\subsection{Can Different Backbone Models Use Scroll Effectively?}
\label{sec:results-backbone}
% ============================================================

Scroll provides an environment for managing context but does not dictate its use: what state to keep where, and what code to write, are left to the model. A natural question is \emph{whether the ability to use Scroll effectively is specific to one backbone or shared across models of varying capability}. We rerun both regimes across six backbones with the harness, tools, prompts, and context-management rules held fixed (Table~\ref{tab:backbone}).

Every backbone can use Scroll, but stronger models benefit more. On LongMemEval$_S$, where the model queries the history database with short programs, all backbones benefit similarly---even the 35B model reaches $88.8$, within six points of the best ($94.8$). On BEAM$_{10M}$ the gap stays within $15$ points. 
On LOCA, however, tasks demand longer trajectories and
more complex program synthesis, so the spread widens to $64$ points at 256K
($86.7$ vs.\ $22.7$). Failures are not protocol-level: all backbones adhere to the CodeAct interface, but weaker models commit more execution errors or terminate prematurely on aggregation-heavy tasks. Scroll's ceiling on such tasks thus rises with multi-step query planning and the ability to decide when evidence suffices, suggesting room for post-training on frontier model traces.


\begin{table}[t]
\centering
\caption{Scroll across backbones. Only the foundation model changes; harness,
tools, prompts, and context-management rules are identical.}
\label{tab:backbone}
\small
\begin{tabular}{lcccc}
\toprule
Backbone & LongMemEval$_S$ & BEAM$_{10M}$ & LOCA 128K & LOCA 256K \\
\midrule
Qwen3.8-Max      & 94.8 & 73.1 &89.3 & 86.7 \\
Qwen3.7-Max      & 92.8 & 66.6 &78.7 & 60.0 \\
Deepseek-v4-pro  & 93.2 & 70.2 &69.3 & 58.7 \\
GLM-5.2          & 93.6 & 70.7 &66.7 & 62.7 \\
Kimi-K2.7        & 92.0 & 67.2 &30.7 & 32.0 \\
Qwen3.6-35B-A3B  & 88.8 & 58.1 & 37.3 & 22.7 \\
\bottomrule
\end{tabular}
\end{table}



% \begin{table}[t]
% \centering
% \caption{Scroll across backbones. Only the foundation model changes; harness,
% tools, prompts, and context-management rules are identical.}
% \label{tab:backbone}
% \small
% \begin{tabular}{lccccc}
% \toprule
% Backbone & LongMemEval$_S$ & BEAM$_{10M}$ (w/o index) & BEAM$_{10M}$ & LOCA 128K & LOCA 256K \\
% \midrule
% Qwen3.8-Max      & 94.8 & 71.3 & 73.1 &89.3 & 86.7 \\
% Qwen3.7-Max      & 92.8 & 70.0 & 66.6 &78.7 & 60.0 \\
% Deepseek-v4-pro  & 93.2 & 67.8 & &69.3 & 58.7 \\
% GLM-5.2          & 93.6 & 67.1 & 70.7 &66.7 & 62.7 \\
% Kimi-K2.7        & 92.0 & 62.8 & &30.7 & 32.0 \\
% Qwen3.6-35B-A3B  & 88.8 & 62.4 & & 37.3 & 22.7 \\
% \bottomrule
% \end{tabular}
% \end{table}




% ============================================================
\subsection{Ablation Study}
\label{sec:ablation}
% ============================================================
\begin{figure}[t]
\centering
\includegraphics[width=0.8\textwidth]{sections/fig/beam_ablation.pdf}
\vspace{-2mm}
\caption{Ablating Scroll's components on BEAM$_{\text{10M}}$ (judge scores; hatched bars are Scroll variants; Qwen3.8-Max, thinking on). Lossy summarization: summaries replace the originals at ingestion. Scroll w/o REPL: \texttt{ms} exposed as ordinary tool calls, with no persistent kernel. Scroll w/o index: the eviction index is removed, leaving keyword search only. Scroll: the full system.}
\label{fig:ablation}
\end{figure}

 
We ablate the core components of Scroll on BEAM$_{10M}$. 
Figure~\ref{fig:ablation} reports judge scores per category and overall.
First, to assess the utility of the Event Log, we compare Scroll against a lossy variant whose history is summarized at ingestion, with the originals discarded. 
Second, to evaluate the programmatic interface, we compare against \emph{Scroll w/o REPL}, which exposes \texttt{search}, \texttt{expand}, and \texttt{sql\_query} as ordinary tool calls, with no persistent kernel. 
Third, to assess index-guided navigation, we remove the eviction index, leaving the agent to locate history through keyword search alone rather than index ranges.


Discarding the original records is the most damaging ablation: the lossy variant falls to 19.9 overall, with near-zero scores wherever the answer must preserve exact values from the history, such as information extraction, temporal reasoning, and knowledge update. 
Scroll w/o REPL underperforms full Scroll by 7.3 points, since serialized tool results cannot be filtered, joined, or aggregated in the kernel; the difference is concentrated in abilities that require composing evidence from many records, such as knowledge update (92.5 vs.\ 82.5) and instruction following (97.5 vs.\ 76.3), while single-lookup abilities are unaffected. 
Removing the eviction index costs 1.8 points overall, but the effect concentrates where evidence is scattered across the history and must otherwise be collected by keyword search: preference following (89.1 vs.\ 74.9), summarization (70.5 vs.\ 62.6), and event ordering (64.1 vs.\ 58.1).

% ============================================================
\subsection{Cost and Efficiency}
\label{sec:results-cost}
% ============================================================
Scroll exposes only a small fraction of the corpus to the model. Ingestion involves no additional LLM calls, and at query time records are filtered inside the Python kernel, so only printed output enters the context. Figure~\ref{fig:cost} shows the per-task distribution of input tokens, output tokens, and agent turns: median input on BEAM$_{10M}$ is $105$K tokens, about $1\%$ of the corpus, and output is an order of magnitude smaller than input across all three benchmarks. Note that for LongMemEval$_S$ and BEAM$_{10M}$ we measure retrieval alone, whereas for LOCA we measure full task completion, hence its longer trajectories. We report token counts rather than latency or dollar cost, as both depend on serving configuration.
 
\begin{figure*}[t]
\centering
\includegraphics[width=0.9\textwidth]{sections/fig/cost.pdf}
\caption{Per-task cost of Scroll (backbone: Qwen3.8-Max): (a) input tokens, (b) output tokens, (c) agent turns. Boxes span the IQR with the median marked; red diamonds are means. Log scale in (a) and (b).}
\vspace{-2mm}
\label{fig:cost}
\end{figure*}
